\documentclass{quiz}
\usepackage{logic}
\usepackage{booktabs}

% Show answers
\noprintanswers % Question Paper
% \printanswers % Solution Paper

% Document info
\title{\textbf{练习题（二）：多值逻辑}}
\author{唯理书院2026\\精确的逻辑与模糊的意义}
\date{}

\begin{document}
\maketitle


\vspace{5ex}
\section*{\centering 符号速查表}

  \begin{center}
  \begin{tabular}{cll}
    \textbf{符号} & \textbf{名称} & \textbf{自然语言对应} \\\midrule
    $\neg$ & 否定 & “不”、“并非” \\
    $\land$ & 合取/逻辑与 & “并且” \\
    $\lor$ & 析取/逻辑或 & “或者” \\
    $\implies$ & 材料蕴含 & “如果...那么...” \\
    $\iff$ & 双条件 & “当且仅当” \\
    $\exists$ & 存在量词 & “存在”、“有些” \\
    $\forall$ & 全称量词 & “所有”、“任意” \\
    $=$ & 等词 & “等于”、“就是” \\
  \end{tabular}
  \end{center}


\vspace{2ex}
\section*{\centering 武卡谢维奇真值表}
\vspace{-1.5ex}

%  \begingroup
%  \setlength{\tabcolsep}{4pt}
  \begin{table}[h]
  \centering
    \begin{tabular}{c|c}
      $\neg$ & \\\hline
      1 & 0 \\
      0 & 1 \\
      \hash & \hash
    \end{tabular}
    \quad
    \begin{tabular}{c|ccc}
      $\land$ & 1 & 0 & \hash \\\hline
      1 & 1 & 0 & \hash \\
      0 & 0 & 0 & 0 \\
      \hash & \hash & 0 & \hash
    \end{tabular}
    \quad
    \begin{tabular}{c|ccc}
      $\lor$ & 1 & 0 & \hash \\\hline
      1 & 1 & 1 & 1 \\
      0 & 1 & 0 & \hash \\
      \hash & 1 & \hash & \hash
    \end{tabular}
    \quad
    \begin{tabular}{c|ccc}
      $\implies$ & 1 & 0 & \hash \\\hline
      1 & 1 & 0 & \hash \\
      0 & 1 & 1 & 1 \\
      \hash & 1 & \hash & 1
    \end{tabular}
  \end{table}
%  \endgroup


\vspace{1ex}
\section*{\centering 克莱尼真值表}
\vspace{-1.5ex}

%  \begingroup
%  \setlength{\tabcolsep}{4pt}
  \begin{table}[h]
  \centering
    \begin{tabular}{c|ccc}
      $\implies$ & 1 & 0 & \hash \\\hline
      1 & 1 & 0 & \hash \\
      0 & 1 & 1 & 1 \\
      \hash & 1 & \hash & \hash
    \end{tabular}
  \end{table}
%  \endgroup


\clearpage
\begin{questions}

  \question \textbf{热身：}用真值表证明下列语句在二元逻辑中是重言式：
  \begin{parts}
    \part $\neg(\phi \land \neg\phi)$
    \part $(\phi \implies \psi) \lor (\psi \implies \phi)$
    \part $\neg\phi \implies (\phi \implies \psi)$
  \end{parts}

\end{questions}



\questionsection{三值逻辑}

在三值逻辑中，任何命题$P$可以取$\set{1, 0, \hash}$中的真值。

\begin{questions}

  \question 用武卡谢维奇的逻辑系统，写出下式完整的真值表（应包含$P, Q$全部$9$种赋值组合）：
    \[ (P \implies Q) \lor (Q \implies P) \]
  
  \vspace{2.5in}
  
  \question 分别使用武卡谢维奇和克莱尼真值表，计算出以下式子在$P = \hash$时的真值。
  \begin{parts}
    \part $P \implies P$
    \part $(P \implies P) \land (P \implies P)$
    \part $P \lor \neg P$
  \end{parts}
  \clearpage
  
  \question 构造一个式子，满足：
  \begin{itemize}
    \item 在武卡谢维奇逻辑中是永真式（所有赋值下真值都为$1$）
    \item 在克莱尼逻辑中不是永真式（至少存在一种赋值使其真值$\neq 1$）
  \end{itemize}
  写出该式子，并分别给出两个完整的真值表来证明其满足条件。
  \vspace{3in}
  
  \question \em{（选做）}克莱尼逻辑中是否存在任何永真式？请给出一个具体例子，或证明不存在。
  
\end{questions}



\questionsection{超赋值逻辑}

记“张三是高的”为$P$。假设所有容许的精细化方式，是把“高”的身高门槛设定为$175$cm到$185$cm之间（含端点）的任意整数厘米值。已知张三身高$179$cm，李四身高$183$cm。

\begin{questions}
  
  \question 在上述精细化集合下，$P$在哪些精细化下为真？哪些为假？$P$是超真、超假，还是既非又非？
  \vspace{2in}
  
  \question 考虑式子$P \lor \neg P$（“张三是高的，或张三不是高的”）。这句式子是超真、超假，还是既非又非？
  \vspace{2in}
  
  \question 记$Q$为“李四是高的”。考虑条件句$P \implies Q$（“如果张三是高的，那么李四也是高的”）。这句式子是超真、超假，还是既非又非？


  
\end{questions}




\end{document}
